% !TEX root = ../main.tex

\chapter{全文总结}

\section{主要结论}

本文针对 WebRTC 实时通信中拥塞控制算法在复杂网络环境下适应性不足的问题，提出了基于离线强化学习的拥塞控制优化与部署验证方案，围绕平台搭建、数据构建、模型训练、推理部署与 A/B 测试评估五个环节形成了完整的端到端技术路线。通过系统性的研究工作，本文得出以下主要结论：

第一，本文构建的 WebRTC 采集部署验证平台证明了利用廉价终端设备即可搭建支撑离线强化学习完整闭环的实验基础设施。平台通过网络轨迹采集、IQL算法、A/B 测试与信道控制的模块化隔离设计，实现了数据采集与策略评测的严格物理解耦，为后续各环节提供了高保真度的底层支撑。

第二，本文设计的数据清洗与数据集构建流程有效解决了真实 WebRTC 采集数据中普遍存在的异常零点、短时缺口与长断裂等质量问题。通过上下文活跃度判定的无效值标记、短缺口线性插值修复与长缺口处的 segment 分段等机制，最终产出了包含 109,393 条有效转移样本、覆盖 1,191 个连续有效轨迹段的高质量训练数据集，为离线策略学习提供了可靠的数据基础。

第三，基于隐式 Q 学习（IQL）算法训练得到的拥塞控制策略在 A/B 测试中展现出显著的动作平滑性优势。在所有九类测试场景下，RL 策略的码率动作波动（cap\_delta）相比 GCC 下降 42\%--73\%，表明策略学到了更稳定的带宽控制行为，能有效抑制 GCC 在弱网反馈噪声下的振荡问题。

第四，离线强化学习策略在多数弱网场景中实现了 RTT 尾部时延与丢包率的同步降低。在 3G、LTE、低带宽、高时延及极端弱网等场景下，RTT P95 下降约 10\%--16\%，丢包率下降约 18\%--25\%，综合 QoE 分数获得正向增益，验证了以少量吞吐让步换取交互稳定性提升的策略设计逻辑。

第五，本文在部署实践中发现，当推理请求与被测媒体流共享物理信道时，远端推理方案在带宽受限场景下存在原理性缺陷——网络拥塞会导致推理延迟膨胀，形成“控制滞后→拥塞加剧→推理更慢”的恶性循环。这一工程发现表明，在实时控制系统中引入深度学习策略时，推理延迟与被控系统的耦合关系必须作为架构设计的首要约束。浏览器本地 ONNX 推理方案通过将控制回路完全封闭于浏览器进程内部，从根本上解决了该问题。

\section{研究展望}

尽管本文已构建了从数据采集到在线部署验证的完整闭环，但仍存在若干值得进一步探索的方向：

第一，奖励函数与策略偏置的优化。当前策略在高质量网络场景中存在带宽利用不足的问题，后续可考虑引入自适应的奖励权重调节机制，或在训练数据中增加更丰富的高带宽利用率专家轨迹，以缓解策略的过度保守倾向。

第二，训练数据规模与多样性的扩展。现有数据集虽已覆盖九类典型网络场景，但在设备类型、浏览器版本、视频内容多样性等方面仍有扩展空间。更大规模、更异构的训练数据有望进一步提升策略的泛化能力。

第三，多模态特征融合。当前状态表示仅包含传输层网络度量，未来可将应用层媒体特征（如视频帧量化参数、关键帧编码大小等）融入状态空间，赋予模型跨层联合优化的能力。

第四，在线微调与持续学习。在浏览器本地 ONNX 推理架构成熟的基础上，可进一步探索联邦学习等隐私保护机制，使终端节点在本地运行策略并收集反馈，周期性向云端聚合梯度，实现拥塞控制模型的持续进化。

第五，更完善的 QoE 评价体系。当前采用的 QoE 代理函数尚未直接覆盖帧率、分辨率变化、卡顿时长等视频主观体验因素，后续可结合视频质量评估模型（如 VMAF）构建更贴近真实用户感知的评估指标。